\setcounter{page}{93}
\section{\S4 The Derived Functors of $\otimes$ and $f^*$}

Let $X$ be a prescheme, and let $F^\bullet,G^\bullet\in K(\operatorname{Mod}(X))$.
Define the tensor product complex $F^\bullet\otimes G^\bullet$ as the simple complex associated to the double complex $(F^p\otimes G^q)$, with
\[
\big(F^\bullet\otimes G^\bullet\big)^n
= \sum_{p+q=n} F^p\otimes G^q,
\]
and differential
\[
d = d_F + (-1)^n d_G.
\]

Homotopies are compatible with tensor products, so we have a bifunctor
\[
\otimes\colon K(\operatorname{Mod}(X)) \times K(\operatorname{Mod}(X))
\to K(\operatorname{Mod}(X)).
\]

\textbf{Lemma 4.1.}
Let $F^\bullet,G^\bullet$ be complexes of $\sheaf{O}_X$-modules, with $G^\bullet$ bounded above.
Assume either
\begin{enumerate}
\item[(a)] $G^\bullet$ is acyclic; or
\item[(b)] $F^\bullet$ is acyclic, and either
        \begin{enumerate}
        \item[1)] $F^\bullet$ is bounded above, or
        \item[2)] $G^\bullet$ is bounded in both directions.
        \end{enumerate}
\end{enumerate}
Then $F^\bullet\otimes G^\bullet$ is acyclic.

\setcounter{page}{94}
\textit{Proof.}
Let $K^{\cdot\cdot}$ denote the double complex $K^{pq}=F^p\otimes G^q$.
We have the standard spectral sequences (EGA $\mathrm{O_{III}}$ 11.3.2):
\[
E_2^{pq} = H_\mathrm{I}^p H_\mathrm{II}^q(K^{\cdot\cdot}) \implies H^n(F^\bullet\otimes G^\bullet),
\]
\[
E_2^{pq} = H_\mathrm{II}^p H_\mathrm{I}^q(K^{\cdot\cdot}) \implies H^n(F^\bullet\otimes G^\bullet).
\]

The hypotheses guarantee biregularity of the spectral sequences.
Case (a): Each $B^q(G^\cdot)=Z^q(G^\cdot)$ is flat, so $F^p\otimes G^q$ is acyclic for all $p,q$.
Hence all $E_2^{pq}=0$, so $H^n(F^\bullet\otimes G^\bullet)=0$.

Case (b): Each $G^q$ is flat, so $F^\bullet\otimes G^q$ is acyclic for all $p,q$.
Again $E_2^{pq}=0$, forcing acyclicity of $F^\bullet\otimes G^\bullet$.

Let $L\subseteq K^-(\operatorname{Mod}(X))$ be the triangulated subcategory of complexes of flat $\sheaf{O}_X$-modules.
By Proposition 1.2 and Lemma 4.1(a), the tensor functor
\[
-\otimes-\colon K^-(\operatorname{Mod}(X)) \to K(\operatorname{Mod}(X))
\]
satisfies the hypotheses of Theorem I.5.1, so we obtain
\setcounter{page}{95}
\[
\mathbf{L}_\mathrm{II}\otimes\colon K^-(\operatorname{Mod}(X)) \times D^-(X) \to D(X).
\]

By Lemma 4.1(b), the functor is exact in the first variable and descends to
\[
\mathbf{L}_\mathrm{I}\mathbf{L}_\mathrm{II}\otimes
\colon D^-(X)\times D^-(X)\to D(X).
\]

Deriving in either order yields the same functor (Lemma I.6.3).
We write simply
\[
F^\bullet\stackrel{\mathbf{L}}{\otimes} G^\bullet
\]
for the left derived tensor bifunctor.

\textbf{Definition.}
For $F^\bullet,G^\bullet\in D^-(X)$, define the \textit{local hyperTor} groups
\[
\mathcal{Tor}_i\big(F^\bullet,G^\bullet\big)
= H^{-i}\big(F^\bullet\stackrel{\mathbf{L}}{\otimes} G^\bullet\big).
\]

\textbf{Proposition 4.2.}
Let $X$ be a prescheme, $F^\bullet\in K^b(\operatorname{Mod}(X))$.
The following are equivalent:
\begin{enumerate}
\item[(i)] There exists a quasi-isomorphism $G^\bullet\to F^\bullet$ with $G^\bullet$ a bounded complex of flat $\sheaf{O}_X$-modules.
\item[(ii)] The functor $F^\bullet\otimes-\colon D^-(X)\to D^-(X)$ is way-out right (hence way-out both ways).
\setcounter{page}{96}
\item[(iii)] There exists $n_0\in\mathbb{Z}$ such that $\mathcal{Tor}_i(F^\bullet,G)=0$ for all $\sheaf{O}_X$-modules $G$ and all $i>n_0$.
\end{enumerate}

\textit{Proof.}
Analogous to Proposition I.7.6; omitted.
Key facts:
1. An $\sheaf{O}_X$-module $F$ is flat $\iff \mathcal{Tor}_1(F,G)=0$ for all $G$.
2. For a surjection $Z\to B\to 0$, one has $Z\otimes G\simeq B\otimes G$ under suitable conditions.
3. Use truncations $\sigma_{\ge n}(F^\bullet)$ to construct bounded flat resolutions.

\textbf{Remark.}
In condition (iii), it suffices to test only quasi-coherent sheaves, or even the stalks $G=\sheaf{O}_{X,x}$ and residue fields $k(x)$ at each $x\in X$.
In (ii), it suffices to restrict the functor to $D_\mathrm{qc}^-(X)$.

\setcounter{page}{97}
\textbf{Definition and Corollary 4.3.}
A complex $F^\bullet\in K^b(\operatorname{Mod}(X))$ satisfying the equivalent conditions of Proposition 4.2 is said to have \textit{finite Tor-dimension}.
These complexes form a localizing subcategory of $K^b(\operatorname{Mod}(X))$, denoted
\[
K^b(\operatorname{Mod}(X))_{\mathrm{fTd}}.
\]
The corresponding derived subcategory is
\[
D^b(X)_{\mathrm{fTd}}.
\]
\textit{Proof.} Cf.\ Corollary I.7.7.

Now consider
\[
-\otimes-\colon K^b(\operatorname{Mod}(X))_{\mathrm{fTd}} \to K(\operatorname{Mod}(X)).
\]
Let $L\subseteq K^b(\operatorname{Mod}(X))_{\mathrm{fTd}}$ be complexes of flat modules.
By Proposition 4.2 and Lemma 4.1(a2), $L$ satisfies Theorem I5.1, so we obtain
\[
\mathbf{L}\otimes\colon K(\operatorname{Mod}(X)) \times D^b(X)_{\mathrm{fTd}} \to D(X).
\]
By Lemma 4.1(b2), this descends to
\[
\mathbf{L}_\mathrm{I}\mathbf{L}\otimes
\colon D(X)\times D^b(X)_{\mathrm{fTd}}\to D(X),
\]
also denoted simply $\stackrel{\mathbf{L}}{\otimes}$.

\setcounter{page}{98}
\textbf{Problem.}
Does the bifunctor
\[
\otimes\colon K(\operatorname{Mod}(X)) \times K^b(\operatorname{Mod}(X))_{\mathrm{fTd}}
\to K(\operatorname{Mod}(X))
\]
admit a left derived functor in the first variable?
For fixed $G^\bullet\in K^b(\operatorname{Mod}(X))_{\mathrm{fTd}}$, does there exist a subcategory $L\subseteq K(\operatorname{Mod}(X))$ satisfying the hypotheses of Theorem I.5.1 for $-\otimes G^\bullet$?

\textbf{Proposition 4.4.}
Let $X$ be a prescheme (resp.\ locally noetherian), $F^\bullet\in D_\mathrm{qc}(X)$ (resp.\ $D_\mathrm{c}(X)$),
$G^\bullet\in D_\mathrm{qc}^-(X)$ (resp.\ $D_\mathrm{c}^-(X)$).
Assume either
\begin{enumerate}
\item[(a)] $F^\bullet\in D^-(X)$, or
\item[(b)] $G^\bullet\in D^b(X)_{\mathrm{fTd}}$.
\end{enumerate}
Then
\[
F^\bullet\stackrel{\mathbf{L}}{\otimes} G^\bullet
\in D_\mathrm{qc}(X)\ (\text{resp. } D_\mathrm{c}(X)).
\]

\textit{Proof.}
Local on $X$. For quasi-coherent/coherent sheaves, $\mathcal{Tor}_i(F,G)$ is quasi-coherent/coherent for all $i$.
Use affine open covers and resolutions by direct sums of $\sheaf{O}_X$ to compute Tor groups; apply Proposition I.7.3.

\setcounter{page}{99}
Now let $f\colon X\to Y$ be a morphism of preschemes.
The pullback functor
\[
f^*\colon \operatorname{Mod}(Y)\to\operatorname{Mod}(X)
\]
admits a left derived functor
\[
\mathbf{L}f^*\colon D^-(Y)\to D^-(X),
\]
since there are enough flat $\sheaf{Y}$-modules, and they are $f^*$-acyclic.

If $f^*$ has finite cohomological dimension on $\operatorname{Mod}(Y)$, we say $f$ has \textit{finite Tor-dimension}, and extend
\[
\mathbf{L}f^*\colon D(Y)\to D(X).
\]

\textbf{Proposition 4.5.}
Let $f\colon X\to Y$ be a morphism of preschemes (resp.\ locally noetherian).
Then $\mathbf{L}f^*$ maps $D_\mathrm{qc}^-(Y)$ into $D_\mathrm{qc}^-(X)$ (resp.\ $D_\mathrm{c}^-(Y)\to D_\mathrm{c}^-(X)$).
If $f$ has finite Tor-dimension, the statement extends to unbounded complexes.

\textit{Proof.} Left to the reader.